\chapter{Sprint 16 --- Mobile: Dashboard 3 màn (Đơn hàng / Công nợ / Hỗ trợ) (2026-06-14)}

\section{Bối cảnh}

BA up frame \textbf{mobile\_dashboard} (\texttt{2474:2}, page Dashboard, file copy
\texttt{aoeiDYlg6vlhJZg2w6Q7o5}) gồm 4 màn cùng thế hệ với Home (Sprint 10).
\texttt{Screen\_1\_Home} đã build, user yêu cầu dựng nốt 3 màn còn lại. Khác Sprint
13/15, frame này đọc được structure qua MCP (text + layout + màu), không phải ảnh
phẳng.

\textbf{Mapping (user chốt):}
\begin{itemize}
  \item \texttt{Screen\_2\_Orders} (\texttt{2474:119}) \(\rightarrow\) bản mobile của
        \texttt{/portal/delivery} (``giao hàng oke'').
  \item \texttt{Screen\_3\_Debt\_Knowledge} (\texttt{2474:233}) \(\rightarrow\) trang
        \textbf{MỚI} \texttt{/portal/debt} (module \texttt{wujia\_portal\_debt}).
  \item \texttt{Screen\_4\_Support} (\texttt{2474:346}) \(\rightarrow\) bản mobile của
        \texttt{/portal/support}.
\end{itemize}

\textbf{Quy tắc user giao:} model/controller chưa có \(\rightarrow\) làm UI trước, note
lại, BA fill backend sau. Frame \texttt{2474} là thế hệ \textbf{CŨ} (header 112px,
footer active ĐỎ) --- shell hiện tại đã FINAL (Sprint 12/14: header cyan 104, footer
active CYAN, sheet ``Thêm''). \textbf{GIỮ SHELL FINAL, chỉ build content} --- footer
đỏ Figma là drift có chủ đích, đối chiếu BA. ADR-007 (mỗi feature portal = 1 module)
\(\rightarrow\) tạo \texttt{wujia\_portal\_debt} (module thứ 19; backend Phase 2 sẽ
vào đây).

\section{Spec theo Figma}

Pattern chung mỗi section: \textbf{title row} (title Inter 18/700 \texttt{\#1F2933} +
``Xem tất cả \(\rangle\)'' 12/700 cyan căn phải) + \textbf{card trắng} border
\texttt{\#E9EEF2} bo \(\sim\)16. Tile icon \(40\times40\) bo 10 nền \texttt{\#E9F7FC}
(\(\approx\) \texttt{--wujia-primary-soft}).

\textbf{Screen 2 --- \texttt{/portal/delivery}:}
\begin{itemize}
  \item \emph{Đơn hàng gần đây}: 1 card gộp 3 row (separator), tile
        \texttt{icon-file-text} + mã 16/700 + giờ ngày 12 + tiền 13/700 + badge soft +
        chevron. Link \texttt{/portal/purchase-history/<id>}.
  \item \emph{Giao hàng sắp tới}: card RỜI mỗi batch, tile \texttt{icon-truck} + tên +
        ngày kiểu ``29/05/2026 (Thứ 5) \textbullet{} 08:00'' + ``Tổng đơn: N đơn'' +
        ``Tổng tiền'' + badge (Đang giao = info / Chuẩn bị giao = muted).
  \item \emph{Yêu cầu đổi trả gần đây}: card rời (không tile) mã + ngày + lý do +
        tổng tiền + badge (Chờ xử lý = danger / Đang xử lý = warning).
\end{itemize}

\textbf{Screen 3 --- \texttt{/portal/debt} (MỚI):}
\begin{itemize}
  \item \emph{Công nợ tóm tắt}: 1 card 3 row kv (tile cyan / tile danger ``Quá hạn'') ---
        Tổng công nợ / Quá hạn thanh toán / Kỳ công nợ hiện tại. \textbf{UI-only}:
        \texttt{account.move} = Phase 2 \(\rightarrow\) ``0 đ'' (precedent KPI Công nợ
        Home Sprint 10) + kỳ ``---'', ``Xem tất cả'' href \texttt{\#}.
  \item \emph{Bài viết / Kiến thức mới}: 1 card 3 row (icon chung \texttt{icon-file-text})
        --- data thật \texttt{wujia.knowledge.article} published, link
        \texttt{/portal/knowledge/<slug>}.
  \item \emph{Đơn hàng gần đây}: y hệt Screen 2 (dùng chung class + helper).
\end{itemize}

\textbf{Screen 4 --- \texttt{/portal/support}:}
\begin{itemize}
  \item \emph{Hỗ trợ nhanh}: 1 card 3 row --- Tạo yêu cầu (\texttt{/portal/support/new},
        thật) / Gọi hotline (\texttt{tel:} từ \texttt{company.phone}) / Chat
        (\textbf{UI-only} href \texttt{\#}, chờ BA channel).
  \item \emph{Thông tin cửa hàng}: 1 card 3 row kv (icon trần 18, không tile) --- Địa chỉ
        / Điện thoại / Người phụ trách (= \texttt{main\_owner\_member\_id.user\_id.name}),
        data thật từ franchise active.
  \item \emph{Giao hàng sắp tới}: y hệt Screen 2.
\end{itemize}

\section{Module \& file chạm}

\begin{itemize}
  \item \texttt{wujia\_portal\_base} (19.0.5.5.0 \(\rightarrow\) 19.0.5.6.0):
    \begin{itemize}
      \item \texttt{controllers/utils.py}: helper CHUNG cho section lặp 2--3 trang ---
            \texttt{get\_recent\_orders(franchise\_ids, limit)},
            \texttt{get\_upcoming\_batches(...)} (trả list dict đã tính sẵn order\_count
            + total \textbf{chỉ của franchise user}), \texttt{format\_batch\_departure}
            (``dd/mm/yyyy (Thứ X) \textbullet{} HH:MM'') + map UI-only
            \texttt{MOBILE\_ORDER\_BADGES} / \texttt{MOBILE\_BATCH\_BADGES}. Guard theo
            \texttt{\_fields} (\texttt{franchise\_id}, \texttt{planned\_departure} không
            thuộc dependency của base) \(\rightarrow\) thiếu module thì trả rỗng, không
            crash.
      \item \texttt{views/portal\_home.xml}: KPI ``Công nợ'' mobile (trước là
            \texttt{<div>} chết) \(\rightarrow\) \texttt{<a href="/portal/debt">}.
    \end{itemize}
  \item \texttt{wujia\_portal\_debt} (MỚI, 19.0.1.0.0): manifest depends
        \texttt{wujia\_portal\_base, wujia\_portal\_knowledge}; controller
        \texttt{/portal/debt} (\texttt{auth='user'}, fetch knowledge + recent\_orders,
        công nợ KHÔNG query); \texttt{views/portal\_debt.xml} 3 section.
  \item \texttt{wujia\_portal\_delivery} (19.0.1.2.0 \(\rightarrow\) 19.0.1.3.0): controller
        \texttt{/portal/delivery} thêm context mobile (recent\_orders / upcoming\_batches
        / recent\_returns + map \texttt{MOBILE\_RETURN\_BADGES}); view bọc desktop
        \texttt{d-none d-lg-block} + block mobile \texttt{d-lg-none .wujia-mpage
        .wujia-mdash} NGOÀI \texttt{.content-wrapper}.
  \item \texttt{wujia\_portal\_support} (19.0.3.2.0 \(\rightarrow\) 19.0.3.3.0): controller
        \texttt{/portal/support} thêm context (\texttt{m\_franchise},
        \texttt{m\_upcoming\_batches}, \texttt{m\_hotline}); view bọc desktop + block
        mobile hub.
  \item \texttt{wujia\_portal\_layout} (19.0.15.0.0 \(\rightarrow\) 19.0.16.0.0):
        \texttt{\_variables.css} 2 token \texttt{--wujia-mdash-tile-bg/-tile-danger-bg};
        \texttt{\_components.css} cụm \texttt{.wujia-mdash-*} (titlerow / card / row gộp +
        separator / card rời / tile / kv / chev) + rule \texttt{.wujia-debt-page} cho
        \(\geq\)992px; \texttt{assets.xml} bump \texttt{\_variables} / \texttt{\_components}
        \(\rightarrow\) \texttt{?v=1124}.
\end{itemize}

\section{Gì đã làm (nghiệp vụ)}

Chủ cửa hàng nhượng quyền trên điện thoại giờ có thêm 3 màn theo đúng thiết kế BA:
trang \textbf{Giao hàng} hiện tóm tắt 3 đơn mới nhất, các chuyến giao sắp tới (kèm số
đơn và tổng tiền \textbf{của riêng cửa hàng mình} trong chuyến) và các yêu cầu đổi trả
gần đây; trang \textbf{Công nợ} mới (bấm từ ô ``Công nợ'' ở Trang chủ) gom số dư công nợ,
bài kiến thức mới và đơn gần đây; trang \textbf{Hỗ trợ} thành một hub có nút tạo yêu cầu,
gọi hotline trực tiếp và xem thông tin liên hệ cửa hàng. Phần số công nợ và chat hiện
là giao diện chờ (account.move thuộc giai đoạn 2), còn lại đều chạy dữ liệu thật.
Bản desktop của cả \texttt{/portal/delivery} và \texttt{/portal/support} GIỮ NGUYÊN
100\%.

\section{Lý do / Trade-off}

\textbf{Helper chung ở \texttt{wujia\_portal\_base}} thay vì t-call cross-module: section
``Đơn hàng gần đây'' + ``Giao hàng sắp tới'' lặp trên 3 trang ở 3 module khác nhau.
Đặt logic query ở base (cả 3 đều depends) + dùng chung \textbf{CSS class}, markup viết
per-view \(\rightarrow\) tránh dependency chéo module (delivery không phải depend support
và ngược lại) mà vẫn DRY phần nặng (query + format). Guard \texttt{\_fields} để base
không hard-depend \texttt{wujia\_sale}/\texttt{wujia\_delivery}.

\textbf{Module mới \texttt{wujia\_portal\_debt}} thay vì nhét route vào module sẵn có:
theo ADR-007, mỗi feature portal = 1 module \(\rightarrow\) khi BA build backend
\texttt{account.move} (Phase 2) thì model + view chi tiết công nợ vào đúng chỗ, không
phình module khác. Trang render \texttt{.wujia-mpage} ở mọi cỡ màn (BA mới vẽ mobile),
\(\geq\)992px cap \texttt{max-width: 720px} qua \texttt{.wujia-debt-page} --- tạm, thay
bằng desktop design thật khi BA vẽ.

\textbf{Badge map MOBILE tách desktop} (precedent \texttt{MOBILE\_STATE\_BADGES} Sprint
13): nhãn theo Figma (``Đã xác nhận'', ``Chờ xử lý''\ldots) lấy từ state thật nhưng map
riêng cho mobile, không đụng map desktop của từng module.

\section{Verify}

\begin{itemize}
  \item \texttt{upgrade.sh} 4 module + install \texttt{wujia\_portal\_debt} RC=0, log
        sạch (ERROR trong log là stale 06-05/06-07, không phải hôm nay).
  \item \texttt{test\_sprint9.py} 7/7 + \texttt{test\_sprint5.py} 0 FAIL (21 PASS ---
        thêm 1 assertion \texttt{S00004 batch\_id} chạy được nhờ batch demo local; không
        phải regression code).
  \item HTTP authenticated (\texttt{anh.owner} local :8019): \texttt{/portal/delivery},
        \texttt{/portal/debt}, \texttt{/portal/support}, \texttt{/portal},
        \texttt{/portal/knowledge}, \texttt{/portal/purchase-history} đều 200; markup
        mobile đủ section mỗi trang; KPI Công nợ Home link \texttt{/portal/debt}.
  \item Batch demo local-only (\texttt{BATCH/OUT/00001}, planned\_departure = mai
        08:00, KHÔNG vào manifest): card render ``Tổng đơn: 1 đơn / Tổng tiền: 299.000 đ''
        \(\rightarrow\) xác nhận lọc franchise đúng (chỉ đơn của user trong batch).
  \item Desktop intact (\texttt{d-none d-lg-block} nguyên vẹn); CSS serve đúng
        \texttt{?v=1124}.
  \item \textbf{Gotcha \#13 lặp lại}: route mới \texttt{/portal/debt} cần restart server
        (\texttt{--dev=reload} không bắt route controller mới).
\end{itemize}

\section{Pending sau Sprint 16}

\begin{itemize}
  \item Công nợ 3 số + ``Xem tất cả'' \(\rightarrow\) wire khi BA build
        \texttt{account.move} (Phase 2).
  \item \texttt{wujia.return.request} CHƯA có field tổng tiền \(\rightarrow\) row đổi trả
        hiện ``---'', wire khi BA thêm field.
  \item Chat channel + giờ chat (UI-only href \texttt{\#}); hotline khi
        \texttt{company.phone} trống.
  \item Mobile CHƯA có đường xem \textbf{list ticket} support (Figma không vẽ) ---
        hỏi BA.
  \item \texttt{/portal/debt} desktop dùng tạm layout mobile cap 720px --- chờ BA design
        desktop.
  \item Có thêm row ``Công nợ'' vào sheet ``Thêm'' không --- hỏi BA.
  \item Pending cũ: trang mobile Đổi trả (\texttt{4449:81}), status giao hàng batch state
        thật (Sprint 13), khung giờ vào cart (Sprint 11), i18n \texttt{.po}.
\end{itemize}
